\documentclass{chsmc}
\chsmcsetup{
	year = 2024,
	part = 1,
	type = B,
	show-answers = false,
}
\begin{document}

\section{填空题：本大题共 8 小题，每小题 8 分，共 64 分.}

\begin{enumerate}
	\item 复数 $z = 1 + \sqrt{3}\ii$（$\ii$ 为虚数单位），则
		$z + 2\overline{z} + 3|z|$ 的模为 \blankline
	\item 设 $a$, $b$ 为实数，若一元二次不等式 $x^2 + b < ax$ 的解集为 $(1,5)$，
		则一元二次不等式 $x^2 - a > bx$ 的解集为 \blankline
	\item 函数 $y = \dfrac{4^x + 1}{2^x + 1}$ 的最小值为 \blankline
	\item 一只青蛙在正方形 $ABCD$ 的四个顶点间跳跃，
		每次跳跃总是等可能地跳至与当前所在顶点相邻的两个顶点之一，
		且各次跳跃是独立的. 若青蛙第一次跳跃前位于顶点 $A$，
		则它第 $6$ 次跳跃后恰好仍位于顶点 $A$ 的概率为 \blankline
	\item 在椭圆 $\Omega$ 中，$F_1$, $F_2$ 为焦点，$A$ 为长轴的一个端点，
		$B$ 为短轴的一个端点，若 $\angle F_1BF_2 = \angle F_1AB$，
		则 $\Omega$ 的离心率为 \blankline
	\item 数列 $\{a_n\}$ 满足：$a_1 = 1$，且对任意正整数 $n$，
		均有 $a_{n+1} = 10^n a_n^2$，则 $\{a_n\}$ 的通项公式为 \blankline
	\item 四面体 $ABCD$ 中，$AB \perp BC$，$BC \perp CD$，
		且 $AB$, $BC$, $CD$, $DA$ 的长分别为 $1$, $2$, $3$, $4$，
		则四面体 $ABCD$ 的体积为 \blankline
	\item 若三个正整数 $a$, $b$, $c$ 的位数之和为 $8$，且组成 $a$, $b$, $c$
		的 $8$ 个数码能排列为 $2$, $0$, $2$, $4$, $0$, $9$, $0$, $8$.
		则称 $(a,b,c)$ 为“幸运数组”，例如 $(9,8,202400)$ 是一个幸运数组.
		满足 $10 < a < b < 100$ 的幸运数组 $(a,b,c)$ 的个数为 \blankline
\end{enumerate}

\section{解答题：本大题共 3 个小题，满分 56 分. 解答应写出文字说明、证明过程或演算步骤.}

\begin{enumerate}[resume]
	\item (本题满分 16 分) 设无穷等比数列 $\{a_n\}$ 的公比 $q$ 满足 $0 < |q| < 1$.
		若 $\{a_n\}$ 的各项和等于 $\{a_n\}$ 各项的平方和，求 $a_2$ 的取值范围.
	\item (本题满分 20 分) 在 $\triangle ABC$ 中，已知
		$\cos C = \dfrac{\sin A + \cos A}{2} = \dfrac{\sin B + \cos B}{2}$，
		求 $\sin C$ 的值.
	\item (本题满分 20 分) 在平面直角坐标系中，将圆心在 $y$ 轴上，
		且与双曲线 $\Gamma\colon x^2-y^2=1$ 的两支各恰有一个公共点的圆称为“好圆”.
		是否存在常数 $\lambda$ 及 $x$ 轴上的一个定点 $A$，满足：
		若两个好圆外切于点 $P$，则它们的圆心距 $d = \lambda|PA|$？
\end{enumerate}

\end{document}
